\begin{tikzpicture}[
  node distance=0.85cm and 0.65cm,
  stage/.style={
    rectangle,
    rounded corners=3pt,
    draw=lc-def,
    fill=lc-def!8,
    line width=0.8pt,
    minimum width=2.25cm,
    minimum height=1.05cm,
    align=center,
    font=\small\bfseries,
    text=lc-ink
  },
  flow/.style={
    lc vector,
    rounded corners=4pt
  },
  feedback/.style={
    lc aux,
    -{Stealth},
    rounded corners=6pt
  }
]
  \node[stage] (q) {提出问题\\{\footnotesize 业务/科学目标}};
  \node[stage, right=of q] (prep) {获取与准备\\{\footnotesize 采集/清洗}};
  \node[stage, right=of prep] (eda) {探索与可视化\\{\footnotesize 规律/假设}};
  \node[stage, right=of eda] (model) {建模与评估\\{\footnotesize 训练/验证}};
  \node[stage, right=of model] (insight) {洞察与决策\\{\footnotesize 预测/解释}};

  \draw[flow] (q) -- (prep);
  \draw[flow] (prep) -- (eda);
  \draw[flow] (eda) -- (model);
  \draw[flow] (model) -- (insight);

  \draw[feedback] (eda.north) -- ++(0,0.55) -| node[pos=0.25, above, lc label, text=lc-note] {数据驱动新问题} (q.north);
  \draw[feedback] (model.south) -- ++(0,-0.55) -| node[pos=0.25, below, lc label, text=lc-note] {特征迭代与数据补充} (prep.south);
  \draw[feedback] (insight.south) -- ++(0,-1.05) -| node[pos=0.25, below, lc label, text=lc-note] {闭环反馈与新周期} (q.south);
\end{tikzpicture}
